\section{Consistency checks at small sizes}\label{app:consistency}

The checks below compare the statements of the paper with exact computations of $T$ for small
$h$. They cannot confirm asymptotic statements, whose error terms are of the same size as the
differences observed, but they would have exposed errors in the finite forms of the bounds. The
polynomials $T$ were computed exactly by an implementation of the definitions in
\cref{sec:pencils} (python-flint rationals, determinants interpolated at $h+1$ rational points;
the script does not import the code that produced the candidate certificates). Each row of
\cref{tab:constants}, with both signs for $a=4,8$, was run for $n=1,2,\dots$ up to $h=32$--$56$; in
addition the numerators $9$ (with $q=35$, $\Ti_2$) and $12$ (with $q=1001$, $\chi_2$), which are
not rows of the theorem, were run to exercise $3\mid a$ with $9\mid a$ and $e=2$ with $a_0=3$.
No finite bound failed in any of the runs (\path{logs/experiments/runs/}).

\subsection{The pole series}
\Cref{lem:pole-series} was tested on five entries $\mu_\xi(\varphi u^s)$ with the series
$\sum_mw^m(\varphi u^s)'(m)$ summed directly (exact rational derivatives, $\xi$ summed with $700$
digits), for the six threshold rows and both signs of $a=4,8$ at $n=1$. The largest relative
discrepancy is $2\cdot10^{-664}$.

\subsection{The prime $2$}
For odd $a$ every entry bound of \cref{prop:two-odd-entries} was checked in the anchored bases
(for example $1936$ entries for $a=5$, $q=829$, $h=44$), with no failure, and the exact
$\vG2(T)$ lies above the finite bound. For $a=4$, $q=995$ and $h=56$, $\vG2(T)=-9156$ against the
finite bound $-9168$; for $a=8$, $q=12061$ the finite bound is attained exactly for every
$h=4n\le48$ ($-63$ at $h=4$, $-10184$ at $h=48$). Since the leading terms are
$C_e(\kappa)h^2-O(h\log h)$ with a logarithmic error of the same size as the differences observed,
we do not tabulate ratios.

\subsection{Odd primes}
For each run and each $n$ the primes dividing $a$, $q$ and $q^2-\sigma a^2$ were tested against the
finite form of the relevant bound: \cref{prop:p-divides-a} (for example for $a=3$, $q=131$, $h=48$,
$\vG3(T)/h^2=-1.2812$ against $-(2\kappa-1)=-1.3333$), \cref{prop:p-divides-q} (the ratio
$\vG p(T)/(fh^2)$ is $0.3332$ for $p=199\mid995$ at $h=56$, and below $1/3$ by amounts
compatible with $O(h\log h)$ for $p=5\mid995$: $0.3084$) and \cref{prop:p-divides-Y}. No bound
failed. The ordinary primes ($p>2K$ and $p^2>2K$ via the jet and pole bounds) were checked
in the earlier numerator-$2$ audit; for the numerators of this paper only the exceptional primes
were tested at each size.

\subsection{Nonvanishing}
\Cref{thm:nonvanishing} predicts $v_p(\Delta(0))=-pt(\gamma+2-t)+t$ at $h=tp$. Exact computation of
$\Delta(0)$ gives $v_{13}=-152$ for $(a,q,\sigma)=(4,995,1)$, $(3,257,-1)$ and $(8,12061,-1)$ with
$h=52$ and $v_{17}=-200$ for $(5,829,1)$ with $h=68$, and $v_{17}=-300$ for $(3,131,1)$ with
$\kappa=7/6$ and $h=102$, in each case as predicted. These are checks of the whole mechanism of
\cref{sec:nonvanishing}, including the reduction of the local functionals, at one prime per row.

\subsection{Size}
The size bound \cref{thm:decay} was not compared with exact determinants for these rows: the
margins are tiny, and the finite sizes available are too small for $O(\log h/h)$ corrections to be
negligible. The comparison of the six exponents with those obtained in the numerator-$2$ work, of
the same order, is the only sanity check.
